\chapter{Knowledge and Skills Evaluation}
\label{chap:knowledge-skills-evaluation}
\raggedbottom

This chapter presents a comprehensive evaluation of the knowledge and skills developed throughout the KU Parking project. It compares the program targets with three complementary perspectives: an AI-based assessment, a student self-evaluation, and advisor feedback. By combining these viewpoints, the chapter aims to provide a more balanced understanding of both technical progress and professional growth. The evaluation is structured around five key competency areas: Software Development; Machine Learning, LLM \& Agentic AI Systems; Data Management and Knowledge Systems; Critical \& Creative Thinking; and Collaboration \& Communication. To support clear interpretation, radar charts are used to visualize the alignment and differences across these competency areas, followed by reflections that explain the reasons behind the observed results.
\flushbottom
\newpage

\section{Comparison Between Program Targets and AI-Based Evaluation}
\label{section:comparison-program-targets-ai-evaluation}

\subsection{Radar Chart Comparing Program Targets and AI Evaluation}
\label{radar-chart-comparing-program-targets-ai-evaluation}
\begin{figure}[H]
    \centering
    \includegraphics[width=0.9\textwidth]{evaluation-charts/ske-ai.png}
    \caption{Radar Chart: Program Targets vs. AI Evaluation}
    \label{fig:ske-ai}
\end{figure}

The radar chart in Figure \ref{fig:ske-ai} compares the SKE expected level with the AI-assessed level across five competency areas: Software Development, Machine Learning, LLM \& Agentic AI Systems, Data Management and Knowledge Systems, Critical \& Creative Thinking, and Collaboration \& Communication. The score is as following:
\begin{itemize}
    \item \textbf{Software Development}: SKE expected level is 2, AI evaluation is 3
    \item \textbf{Machine Learning, LLM \& Agentic AI Systems}: SKE expected level is 3, AI evaluation is 3
    \item \textbf{Data Management and Knowledge Systems}: SKE expected level is 2, AI evaluation is 2
    \item \textbf{Critical \& Creative Thinking}: SKE expected level is 3, AI evaluation is 3
    \item \textbf{Collaboration \& Communication}: SKE expected level is 3, AI evaluation is 3
\end{itemize}

\subsection{Reflection on the AI Evaluation}
\label{reflection-ai-evaluation}
For each competency area, our team reflected on the AI evaluation results as follows:
\begin{itemize}[itemsep=0.2em, topsep=0.2em, parsep=0pt, partopsep=0pt]
    \item \textbf{Software Development}: The AI assessed this area above the expected level, and we agree with this result. Based on evaluator feedback from the project defenses, we refactored the system into a more modular structure by separating the mobile application, web portal, and AI component, and we also introduced caching for vacancy status. These improvements directly strengthened scalability, maintainability, and system responsiveness.
    \item \textbf{Machine Learning, LLM \& Agentic AI Systems}: The AI evaluation matched the expected level. Our team designed and evaluated the AI component using clear model metrics (Precision, Recall, F1 score, mAP50, and mAP50-95), and we compared multiple training-data configurations before selecting the final model. This reflects a solid ability to apply machine learning workflows within project constraints.
    \item \textbf{Data Management and Knowledge Systems}: The AI evaluation matched the expected level, which aligns with our implementation progress. We established data handling mechanisms such as storing human-drawn parking boundaries for future retraining and using caching for frequently requested vacancy information. These practices demonstrate practical data management for both system performance and long-term model improvement.
    \item \textbf{Critical \& Creative Thinking}: The AI evaluation matched the expected level, and we consider this appropriate. Throughout development, we addressed technical trade-offs, such as balancing dataset diversity and model generalization, and adapted our design after receiving feedback from project defenses. These decisions required both analytical reasoning and creative problem-solving.
    \item \textbf{Collaboration \& Communication}: The AI evaluation matched the expected level and reflects our team process. We worked with clear role separation, maintained regular weekly synchronization, and communicated continuously to resolve blockers quickly. In addition, presenting the project during the SKE Project Expo helped us explain technical decisions clearly to advisors, faculty, and peers.
\end{itemize}

\section{Self-Evaluation by the Student}
\label{section:self-evaluation-student}

\subsection{Radar Chart Comparing Program Targets, AI Evaluation, and Self-Evaluation}
\label{radar-chart-comparing-program-targets-ai-evaluation-self-evaluation}
\begin{figure}[H]
    \centering
    \includegraphics[width=0.9\textwidth]{evaluation-charts/ske-ai-student.png}
    \caption{Radar Chart: Program Targets vs. AI Evaluation vs. Student Self-Evaluation}
    \label{fig:ske-ai-student}
\end{figure}

The radar chart in Figure \ref{fig:ske-ai-student} compares the SKE expected level, AI-assessed level, and student self-evaluated level across five competency areas: Software Development, Machine Learning, LLM \& Agentic AI Systems, Data Management and Knowledge Systems, Critical \& Creative Thinking, and Collaboration \& Communication. The score is as following:
\begin{itemize}
    \item \textbf{Software Development}: SKE expected level is 2, AI evaluation is 3, student self-evaluation is 3
    \item \textbf{Machine Learning, LLM \& Agentic AI Systems}: SKE expected level is 3, AI evaluation is 3, student self-evaluation is 3
    \item \textbf{Data Management and Knowledge Systems}: SKE expected level is 2, AI evaluation is 2, student self-evaluation is 2
    \item \textbf{Critical \& Creative Thinking}: SKE expected level is 3, AI evaluation is 3, student self-evaluation is 3
    \item \textbf{Collaboration \& Communication}: SKE expected level is 3, AI evaluation is 3, student self-evaluation is 3
\end{itemize}

\subsection{Reflection on Self-Evaluation Compared with AI Evaluation}
\label{reflection-self-evaluation-compared-ai-evaluation}

For each competency area, we reflected on the AI evaluation results and found that we agree with the AI assessment as follows:
\begin{itemize}[itemsep=0.2em, topsep=0.2em, parsep=0pt, partopsep=0pt]
    \item \textbf{Software Development}: We agree with the AI evaluation in this area, where the assessed level is above the expected target. This is consistent with our practical progress, including improved modularity across the mobile application, web portal, and AI component, as well as better responsiveness through caching.
    \item \textbf{Machine Learning, LLM \& Agentic AI Systems}: We agree with the AI evaluation at the expected level. This matches our development process, where we performed model training, metric-based performance assessment, and iterative dataset experimentation.
    \item \textbf{Data Management and Knowledge Systems}: We agree with the AI evaluation at the expected level and consider it appropriate. Our implementation supports this result through structured storage of human-drawn parking boundaries and efficient vacancy information retrieval.
    \item \textbf{Critical \& Creative Thinking}: We agree with the AI evaluation at the expected level. During development, we continuously applied analytical and creative thinking to handle technical trade-offs and iteratively improved the model based on our hypotheses.
    \item \textbf{Collaboration \& Communication}: We agree with the AI evaluation at the expected level. This reflects our teamwork practice, including clear task ownership, regular communication, and coordinated presentations during project milestones and the SKE Project Expo.
\end{itemize}


\section{Advisor Evaluation}
\label{section:advisor-evaluation}

\subsection{Radar Chart Comparing Program Targets, AI Evaluation, Self-Evaluation, and Advisor Evaluation}
\label{radar-chart-comparing-program-targets-ai-evaluation-self-evaluation-advisor-evaluation}
\begin{figure}[H]
    \centering
    \includegraphics[width=0.9\textwidth]{evaluation-charts/ske-ai-student-advisor.png}
    \caption{Radar Chart: Program Targets vs. AI Evaluation vs. Student Self-Evaluation vs. Advisor Evaluation}
    \label{fig:ske-ai-student-advisor}
\end{figure}

The radar chart in Figure \ref{fig:ske-ai-student-advisor} provides the final comprehensive comparison of the project's performance, incorporating the SKE expected level,
the AI-assessed level, the student self-evaluated level, and the evaluation from the project advisor. The scores across the five competency areas are as follows:
\begin{itemize}[itemsep=0.2em, topsep=0.2em, parsep=0pt, partopsep=0pt]
    \item \textbf{Software Development}: SKE expected level is 2, AI evaluation is 3, student self-evaluation is 3, and advisor evaluation is 3.
    \item \textbf{Machine Learning, LLM \& Agentic AI Systems}: SKE expected level is 3, AI evaluation is 3, student self-evaluation is 3, and advisor evaluation is 1.
    \item \textbf{Data Management and Knowledge Systems}: SKE expected level is 2, AI evaluation is 2, student self-evaluation is 2, and advisor evaluation is 3.
    \item \textbf{Critical \& Creative Thinking}: SKE expected level is 3, AI evaluation is 3, student self-evaluation is 3, and advisor evaluation is 3.
    \item \textbf{Collaboration \& Communication}: SKE expected level is 3, AI evaluation is 3, student self-evaluation is 3, and advisor evaluation is 3.
\end{itemize}

\subsection{Advisor's Comments on the Evaluation}
\label{advisor-comments-evaluation}
The final evaluation and scores were provided by the project advisor. The following is a comment from the advisor regarding the evaluation on each competency area:
\begin{itemize}[itemsep=0.2em, topsep=0.2em, parsep=0pt, partopsep=0pt]
    \item \textbf{Software Development}: The developers have successfully applied a modular design in their project. Scalability has been considered, particularly through the use of caching for computationally expensive operations. Both unit tests and integration tests have been implemented across all major components.
    \item \textbf{Machine Learning, LLM \& Agentic AI Systems}: AI/ML techniques have been incorporated for image processing and automatic recognition of parking slots. However, large language model (LLM) functionality has not yet been implemented.
    \item \textbf{Data Management and Knowledge Systems}: Multiple storage technologies, including \textit{Amazon ElastiCache}, \textit{PostgreSQL}, and \textit{Amazon S3}, have been effectively integrated into the system architecture to address different storage requirements.
    \item \textbf{Critical \& Creative Thinking}: A range of techniques and tools has been employed to develop an innovative system. However, the solution does not yet reach the level of a disruptive technology.
    \item \textbf{Collaboration \& Communication}: The project documentation and presentation slides are well-prepared, featuring a comprehensive set of diagrams, screenshots, and results, accompanied by meaningful analysis and insights.
\end{itemize}

\subsection{Student Reflection on the Advisor's Evaluation}
\label{student-reflection-advisor-evaluation}
For each competency area, our team reflected on the advisor's evaluation results and found that we agree with the assessment as follows:
\begin{itemize}[itemsep=0.2em, topsep=0.2em, parsep=0pt, partopsep=0pt]
    \item \textbf{Software Development}: We agree with the advisor's assessment and are encouraged by the recognition of our architectural decisions regarding modularity. The implementation of a modular cloud-based architecture and the introduction of caching for vacancy status were core focuses of our iterative development process to improve system responsiveness and scalability.
    \item \textbf{Machine Learning, LLM \& Agentic AI Systems}: We accept the score of 1 based on the specific rubric criteria leaning towards Large Language Model (LLM) integration. However, we would like to reflect that our project's primary technical challenge was centered on Computer Vision and object detection using the \textit{YOLO} model. While the current rubric emphasizes LLM-based systems, we believe our work demonstrates a high level of competency in this alternative domain of Machine Learning.
    \item \textbf{Data Management and Knowledge Systems}: We agree with the score of 3 and appreciate the advisor's recognition of our multi-layered storage strategy. Effectively integrating technologies like \textit{Amazon ElastiCache} for high-speed retrieval, \textit{PostgreSQL} for relational data, and \textit{Amazon S3} for archiving retraining data was instrumental in meeting our system goals.
    \item \textbf{Critical \& Creative Thinking}: We agree with the score of 3. Throughout the project, we applied analytical reasoning to handle technical trade-offs, such as balancing dataset diversity and model generalization. We acknowledge that while our solution is innovative for campus navigation, there remains room for further "disruptive" enhancements, such as automated boundary mapping.
    \item \textbf{Collaboration \& Communication}: We agree with the assessment and believe it reflects our rigorous team process. Maintaining regular weekly synchronization and producing comprehensive documentation—including detailed diagrams and performance analysis—was essential for our successful project execution.
\end{itemize}